\documentclass[11pt]{article}

\usepackage[a4paper,margin=1in]{geometry}
\usepackage[T1]{fontenc}
\usepackage[utf8]{inputenc}
\usepackage{lmodern}
\usepackage{hyperref}

\title{Rehab-FlowMatch: Physics-Constrained Flow Matching for Rehabilitation Motion Trajectories}
\author{Arnav Salkade}
\date{\today}

\begin{document}
\maketitle

\begin{abstract}
Rehab-FlowMatch studies whether noisy wearable IMU data from rehabilitation exercises can be turned into clean, physically plausible motion trajectories. The model learns a time-conditioned velocity field that gradually moves a noisy trajectory toward a clean one, while regularizers discourage jerky motion and unsafe joint limits. We evaluate both reconstruction quality and safety-oriented measures such as constraint violations and temporal stability. A secondary track explores a hybrid ANN-SNN version for lower-power deployment and lightweight on-device adaptation.
\end{abstract}

\section{Introduction}
Stroke and musculoskeletal rehabilitation require frequent, consistent movement practice beyond supervised clinical sessions. In home settings, patients often receive limited feedback, and wearable sensors introduce noise, drift, and missing data. This creates a central challenge: how to recover physically valid trajectories from imperfect measurements while maintaining reliability for safety-critical feedback.

Generative trajectory correction methods are promising because they model denoising as a continuous process instead of a single-step mapping. In this project, we use flow matching as the core formulation. Given a noisy trajectory state and a time variable, the model predicts a velocity field that guides the state toward a clean trajectory. This design supports controllable reverse integration and allows constraints to be incorporated during both training and sampling.

Low reconstruction loss alone is not enough for rehabilitation use. A model can score well numerically while producing implausible jerk spikes, unstable timing, or joint-limit violations. Therefore, Rehab-FlowMatch explicitly targets safety-relevant physical validity alongside denoising quality. The project objective is to build a reproducible pipeline that generalizes across patients and exercises while preserving biomechanical realism.

To make the report easier to scan, the sections below move from current status to model direction, then to the main research questions and experimental plan.

\subsection*{Project Objectives}
\begin{itemize}
  \item Learn a stable time-conditioned velocity field for rehabilitation trajectory denoising.
  \item Enforce kinematic plausibility through smoothness and constraint-aware objectives.
  \item Benchmark quality-safety tradeoffs using trajectory and constraint metrics.
  \item Establish reproducible dataset and integrity checks for reliable experimentation.
  \item Evaluate whether hybrid neuromorphic extensions can reduce compute cost without compromising motion safety.
\end{itemize}

\section{Current Project Status}
\subsection*{What Is Working}
The core pipeline is working end to end: the repo loads the processed StrokeRehab data, passes integrity checks, and runs on Apple silicon with MPS support.

\subsection*{Where Training Stands}
The \texttt{hybrid\_mic\_flow} experiment in \path{experiments/hybrid_mic_flow_mps_fullpower} resumed from \path{experiments/hybrid_mic_flow_mps_fullpower/best.pt} and is still training. The best checkpoint so far is epoch 35 with validation masked RMSE 0.3339. The live run has reached epoch 37, step 15100 of 15585.

\subsection*{Latest Validation}
The latest completed validation at epoch 36 reported train loss 0.9900, masked RMSE 0.3399, frame accuracy 0.2680, and macro F1 0.0318.

\subsection*{Takeaway}
The project is past setup and now in the tuning phase: the main question is whether later epochs beat the current best checkpoint without hurting restoration quality.

\section{Model Direction}
Flow matching and rectified-flow style formulations let us treat denoising as a guided path from noisy motion to clean motion without turning the whole problem into a single-step regression. For this project, the model predicts how each motion state should move over time, and the constraint terms keep that motion smooth and biomechanically plausible.

The hybrid ANN--SNN version is a systems-level extension rather than a replacement for the main model. ANN layers handle the precise parts of the predictor, while SNN blocks may provide sparse feature extraction and lower-power operation. Any bridge between the two must be stable and explicit.

\section{High-Value Research Questions}
The questions below are grouped by safety, hybrid design, and deployment/adaptation so the reader can follow the logic from model behavior to system use.

\subsection*{Quality and Safety}
\begin{enumerate}
  \item Can a hybrid ANN--SNN flow field match ANN trajectory quality while reducing latency and energy?
  \item Which constraint strategy is best in hybrid generators: soft penalties, hard projection, or staged combinations?
  \item Does quantization plus hybrid spiking inference preserve kinematic validity, not only RMSE?
\end{enumerate}

\subsection*{Hybrid Design}
\begin{enumerate}
  \item Which layers should remain ANN versus SNN to preserve reverse-sampling stability?
  \item What tradeoff appears between spike sparsity and trajectory smoothness, especially jerk profiles?
  \item Can sampling-time uncertainty flag unsafe trajectories early enough for fallback control?
\end{enumerate}

\subsection*{Deployment and Adaptation}
\begin{enumerate}
  \item Does three-factor STDP, using global modulation from residual or constraint violation, improve patient-specific adaptation?
  \item Does STDP adaptation improve robustness under IMU drift and dropout without full retraining?
  \item How does eventized IMU input compare with dense IMU input for both error and safety metrics?
  \item What edge-cloud deployment split, with an edge SNN safety filter and cloud ANN generator, gives the best real-time feedback performance?
\end{enumerate}

\section{Experimental Priorities}
\subsection*{Priority A: Safety-First Baseline}
First, establish a strong ANN physics-constrained flow-matching baseline with grouped patient splits and explicit safety metrics: trajectory RMSE, joint-limit violations, and jerk-based smoothness. This is the anchor for all hybrid comparisons.

\subsection*{Priority B: Hybrid Partition Study}
Evaluate several ANN--SNN partitions (early, middle, and late SNN blocks) while keeping the time embedding and output head dense. Compare quality, stability, latency, and energy proxies at matched solver settings.

\subsection*{Priority C: Online Adaptation}
Test bounded STDP adaptation in a small adaptor subnetwork under realistic sensor corruption (bias random walk, dropout windows). Report gains and any safety regressions.

\subsection*{Priority D: Deployment Path}
Prototype an edge-cloud split: edge model performs eventization and safety screening; cloud model handles heavier long-horizon generation and clinician analytics. Measure end-to-end loop latency and fallback reliability.

\section{Working Hypothesis}
The most likely high-impact architecture is a constrained flow-matching backbone with selective hybridization: ANN for stability-critical components, SNN for sparse temporal feature handling, and limited online STDP adaptation for patient and sensor drift personalization. The target outcome is not just lower error, but safer and faster rehabilitation feedback under real-world sensing noise.

\section*{Citation Note}
This draft section consolidates the current literature scan. Formal BibTeX references should be added next in Overleaf for submission-ready formatting.

\end{document}
